%%%BLOCK id=092-02 ch=08 sec=简谐运动·判断（小量近似） kind=derivation alt=09 cont=none y=0.265-0.405
扰动后做简谐运动
%FIG p092_b 0.09 0.302 0.28 0.333
\notefig[0.25]{figures/p092_b.png}
\begin{align*}
F_{\text{回}}&=\frac{kQq}{(L-x)^{2}}-\frac{kQq}{(L+x)^{2}}
=\frac{kQq}{L^{2}\left(1-\dfrac{2x}{L}+\underset{\approx0}{\left(\dfrac{x}{L}\right)^{2}}\right)}\\
&=\frac{kQq}{L^{2}}\cdot\frac{\dfrac{4x}{L}}{1-\underset{\approx0}{\dfrac{4x^{2}}{L^{2}}}}=\frac{4kQqx}{L^{3}}\ \propto x
\end{align*}
% CHECK: 中间两个等号处的式子手写较乱（第二项 -kQq/(L+x)^2 似乎未体现），原样照录
故为简谐运动

or\quad $F=\dfrac{4kQqLx}{\left[L^{2}\left(1-\underset{\approx0}{\dfrac{x^{2}}{L^{2}}}\right)\right]^{2}}=\dfrac{4kQq}{L^{3}}\,x$
%%%END
%%%BLOCK id=100-01 ch=08 sec=简谐运动·两种模型 kind=summary alt=none cont=none y=0.09-0.28
% 原稿：左侧一个大括号并列“简谐运动”与“简谐运动的两种模型”；“两种模型”右侧再用大括号括出下面两项
\begin{itemize}
\item 简谐运动
\item 简谐运动的两种模型
  \begin{itemize}
  \item 弹簧振子（水平）
  \item 单摆
  \end{itemize}
\end{itemize}
自由振动\quad $f=f_{\text{固有}}$\quad $T=T_{\text{固有}}$
% “自由振动”一行位于右大括号内、“弹簧振子（水平）”与“单摆”之间偏右处
%%%END
%%%BLOCK id=100-02 ch=08 sec=简谐运动·公式与图像 kind=summary alt=none cont=none y=0.24-0.34
% 以下各项（至“单摆及周期公式”）在原稿中被一个大括号归并在左侧“机械振动”之下
\textbf{机械振动。}
\begin{itemize}
\item 简谐运动的公式和图像。\quad $F_{\text{回}}=-kx$（回复力\ 效果力）指向平衡位置。
  \begin{itemize}
  \item $\uparrow$正余弦曲线。\quad $T=2\pi\sqrt{\dfrac{L}{g}}$\quad $T=2\pi\sqrt{\dfrac{m}{k}}$
  \end{itemize}
\end{itemize}
% “正余弦曲线”上方有向上箭头指向“图像”
%%%END
%%%BLOCK id=100-03 ch=08 sec=受迫振动与共振 kind=summary alt=none cont=none y=0.27-0.43
\begin{itemize}
\item 受迫振动与共振\quad 机械能不守恒
  \begin{itemize}
  \item 系统在驱动力下的振动。
  \item 受迫振动频率等于驱动力频率，当驱动力频率等于系统固有频率时，发生共振。
  \end{itemize}
\end{itemize}
% 原稿：“机械能不守恒”写在“受迫振动与共振”上方偏右处；其左侧另有两行小字如下
机械工作底座振动\\
共振筛、\unclear{声呐}

（共振曲线：纵轴$A$（受迫振动振幅），横轴$f_{\text{驱}}$，峰值对应$f_{\text{固}}$。）
%FIG p100_a 0.775 0.337 0.955 0.406
\notefig[0.28]{figures/p100_a.png}
%%%END
%%%BLOCK id=100-04 ch=08 sec=简谐运动·规律与特征 kind=summary alt=none cont=none y=0.38-0.50
\begin{itemize}
\item 简谐运动的规律。
  \begin{itemize}
  \item 受力特征\quad $F=-kx$
  \item 运动特征：$a$ $F$ $x$ 同\unclear{增}减，$v$ 与其相反
  \item 能量特征
    \begin{itemize}
    \item 振幅越大能量越大，
    \item 动能势能相互转化，机械能守恒
    \end{itemize}
  \item 对称性特征
    \begin{itemize}
    \item 关于$O$对称，$aFx$大小相等的\unclear{相应}，$V$大小相等，方向可同可反，动能势能相等
    \end{itemize}
  \item 周期性特征
    \begin{itemize}
    \item 标量为$\dfrac{1}{2}$\unclear{$T$}，矢量为$T$
    \end{itemize}
  \end{itemize}
\item 简谐运动的理解及应用
  \begin{itemize}
  \item 图像的信息利用
  \end{itemize}
\end{itemize}
%%%END
%%%BLOCK id=100-05 ch=08 sec=单摆·周期公式 kind=formula alt=none cont=none y=0.50-0.58
\begin{itemize}
\item 单摆及周期公式。\quad $F_{\text{回}}$\ $F_{\text{向}}$\ 大小方向时刻变。\quad $T=2\pi\sqrt{\dfrac{L_{\text{等效摆长}}}{g_{\text{等效重力加速度}}}}$
\end{itemize}
\[ F_{\text{回}}=-mg\sin\theta=-mg\frac{x}{L}=-kx\qquad F_{\text{向}}=T-mg\cos\theta \]
%%%END
%%%BLOCK id=100-06 ch=08 sec=机械波·基本概念 kind=summary alt=none cont=none y=0.57-0.77
% 以下各项在原稿中被一个大括号归并在左侧“机械波”之下；左侧三行标注如下
机械波\\
传播形式和能量\\
不传播质点。
\begin{itemize}
\item 机械波。横波和纵波\ $\leftarrow$\ 振动方向与传播方向平行，有密部和疏部。
\item 波源$+$介质\quad $\hookrightarrow$\ 振动方向与传播方向垂直
\item 横波的图像。波长、波速和频率的关系。\quad $v=\lambda f$\quad 同介质，$v$ 相同。
\item 介质中各点 $T$、$A$、$f$ 都与波源相同。一个周期路程 $4A$，位移为0。\quad $\lambda=vT$
\item 波的干涉和衍射，多普勒效应\ 波源与观察者间有相对运动，波源与波\unclear{频}率不变，观察者接收到的频率\unclear{增大}
  \begin{itemize}
  \item $\hookrightarrow$\ 两波频率必须相等。
  \item $\hookrightarrow$\ 波长与障碍物或孔的尺寸差不多或大。
  \end{itemize}
\end{itemize}
%%%END
%%%BLOCK id=100-07 ch=08 sec=波的图像与多解问题 kind=summary alt=none cont=none y=0.77-1.0
\begin{itemize}
\item 机械波的传播与波的图像。
  \begin{itemize}
  \item $y$-$x$ 图
  \end{itemize}
\item 振动图像与波的图像综合应用
  \begin{itemize}
  \item 用时间联系
  \end{itemize}
\item 波的多解问题。
  \begin{itemize}
  \item 双向性
    \begin{itemize}
    \item 传播方向双向性
    \item 振动方向双向性
    \end{itemize}
  \item 周期性
    \begin{itemize}
    \item $\Delta t$ 与 $T$ 关系不确定
    \item $\Delta x$ 与波长关系不确定。
    \end{itemize}
  \item 隐含性\quad 是只给出波形的一部分或只给出几个特殊点。
  \end{itemize}
\item 波的干涉、衍射和多普勒效应
\end{itemize}
%%%END
%%%BLOCK id=102-01 ch=08 sec=用单摆测重力加速度 kind=method exp=1 alt=none cont=none y=0.03-0.58
\textbf{探究单摆的运动\quad 用单摆测定重力加速度}

\textbf{实验原理与操作}
\[ T=2\pi\sqrt{\frac{L}{g}}\qquad\qquad g=\frac{4\pi^{2}L}{T^{2}} \]

\textbf{数据处理}
\begin{itemize}
\item 公式法
\item 图像法.
\end{itemize}

\textbf{误差分析}
\begin{itemize}
\item 偶然误差.
 \begin{itemize}
 \item 测量时间及摆长时产生的误差
 \item 减小方法：多次测量求平均值/从单摆经过平衡位置时开始计时
 \end{itemize}
\item 系统误差
 \begin{itemize}
 \item 单摆自身
 \item 摆球选体积小，密度大的
 \item 最大摆角要小于$5^\circ$.
 \end{itemize}
\end{itemize}

\textbf{注意事项}
\begin{itemize}
\item 用$1\unit{m}$左右细线
\item 悬线顶端不能晃动用夹子夹住，保证顶点固定
\item 在同一竖直面内摆动，摆角小于$5^\circ$. 从摆球摆到平衡位置开始计时，并数准全振动的次数.
\end{itemize}
%%%END
%%%BLOCK id=107-01 ch=08 sec=振动图像与波动图像 kind=summary alt=none cont=none y=0.03-0.27
\textbf{振动}　$y$-$t$　单个质点不同时刻\\
\hspace*{1em}矢量　周期均为$T$.　能量　周期 $\dfrac{T}{2}$\\
\hspace*{1em}$v$和$E_k$看斜率　斜率为正速度正　$a\ v\ x$ 图像\\
\hspace*{1em}势能和$a$看位移　位移为正加速度\unclear{负}　用求导法画\\
\hspace*{1em}应代入\unclear{数据}建立\unclear{方程}
\[
\left\{\begin{aligned}
y&=A\sin(\omega x+\varphi)\\
y&=A\sin(\omega t+\varphi)
\end{aligned}\right.
\qquad
\begin{aligned}
&y\text{-}x\\
&y\text{-}t
\end{aligned}
\]
\textbf{波动}　$y$-$x$　整体一刻图象
%FIG p107_c 0.775 0.137 0.84 0.165
\notefig[0.25]{figures/p107_c.png}

\hspace*{1em}振动方向看传播：同侧法　　起振方向看最后质点起振方\unclear{向}\\
\hspace*{1em}加速度方向看位移　　两点状态看$y$-$t$

\medskip
\textbf{①}　$y$-$x$的某质点对应$y$-$t$某时刻
%FIG p107_a 0.08 0.19 0.44 0.269
\notefig[0.4]{figures/p107_a.png}

峰对峰　谷对谷　平衡对平衡　　$y$-$t$向右看　$y$-$x$顺\unclear{传播}\\
上升对下降　下降对上升　　　1对$b$
%%%END
%%%BLOCK id=107-02 ch=08 sec=波的多解问题 kind=summary alt=none cont=none y=0.26-0.36
\textbf{②}　多解.　左右两种可能\\
空间多解（方向）　时间多解（$n$的选取）
%FIG p107_b 0.355 0.268 0.485 0.348
\notefig[0.3]{figures/p107_b.png}
\[
\Delta X=\left(n+\dfrac{3}{4}\right)\lambda\quad\text{or}\quad \Delta X=\left(n+\dfrac{1}{4}\right)\lambda
\]
\textbf{③}　先传到后振动.
%%%END
%%%BLOCK id=107-03 ch=08 sec=振动与波动解题要点 kind=formula alt=none cont=none y=0.36-0.41
路程　$S=\dfrac{t}{T}\cdot 4A$　$\big|$　\underline{位移}用三角对应法　$\big|$　三角函数找　$\big|$　考虑相长相消（同频波）（波的干涉）$\big|$\\[0.8ex]
同向　$t=\dfrac{S}{v_1+v_2}$　　反向　$t=\dfrac{S}{|v_1-v_2|}$ % CHECK: 原稿“同向”用v1+v2、“反向”用|v1-v2|，与常识(相向相加、同向相减)相反，照原样转录
%%%END
%%%BLOCK id=107-04 ch=08 sec=波速与频率·机械波与电磁波 kind=knowledge alt=14,15 cont=none y=0.40-0.47
\fbox{机械波}　\underline{$V$只与介质有关}　$\big|$　本子中波速全为$V$（不管任何机械波）\\[1ex]
\fbox{电磁波}　$f$不变　$\big|$　$V=\dfrac{c}{n}$、颜色不变
%%%END
%%%BLOCK id=107-05 ch=08 sec=波动方程 kind=formula alt=none cont=none y=0.47-0.54
\[
\text{波动方程}\ \left\{\begin{aligned}
\text{向右}\quad y&=A\sin\left[\omega\left(t-\dfrac{x}{v}\right)+\varphi\right]\\
\text{向左}\quad y&=A\sin\left[\omega\left(t+\dfrac{x}{v}\right)+\varphi\right]
\end{aligned}\right.
\qquad
\begin{aligned}
&y\text{-}t:\ x\text{定}\\
&y\text{-}x:\ t\text{定}
\end{aligned}
\]
%%%END
%%%BLOCK id=107-06 ch=08 sec=波的干涉 kind=knowledge alt=15 cont=none y=0.54-0.66
\fbox{波的干涉}　波程差为
\[
\left\{\begin{aligned}
2n\cdot\dfrac{\lambda}{2}&\quad\text{相长}\quad A_1+A_2\\
(2n+1)\cdot\dfrac{\lambda}{2}&\quad\text{相消}\quad |A_1-A_2|
\end{aligned}\right.
\]
前提同频波　$f_1=f_2$\\
看起振　振动方向，步调始终相反\\
若为反相波则　规律相反　　半波损失$\to$反相波\\[1ex]
$\Delta X=2n\cdot\dfrac{\lambda}{2}$　$n=0$为最亮纹　$n=1$为第一亮纹…\\[1ex]
$\Delta X=(2n+1)\dfrac{\lambda}{2}$　第零\unclear{亮}纹 % CHECK: (2n+1)λ/2 应对应暗纹，此处字迹像“亮”，存疑
%%%END
%%%BLOCK id=107-07 ch=08 sec=波的衍射 kind=knowledge alt=15 cont=none y=0.63-0.67
\fbox{波的衍射}　$\lambda\ge$ 或 \unclear{$\approx$} 障碍物尺寸，出现\underline{明暗相间亮条纹}
%%%END
%%%BLOCK id=107-08 ch=08 sec=受迫振动与共振 kind=knowledge alt=none cont=none y=0.67-0.78
\[
\left\{\begin{aligned}
&\text{受迫振动}\quad f=f_{\text{力}}\quad\text{共振}\ \bigl(f_{\text{力}}=f_{\text{固有}}\text{时，能量最大}\bigr)\\
&\text{自由振动}\quad f=f_{\text{固有}}
\end{aligned}\right.
\]
%FIG p107_d 0.58 0.672 0.745 0.775
\notefig[0.3]{figures/p107_d.png}
%%%END
%%%BLOCK id=108-01 ch=08 sec=单摆实验 kind=knowledge exp=1 alt=none cont=none y=0.03-0.18
\textbf{单摆}\quad\textbf{①}\ 器材
\begin{itemize}
\item 球：\unclear{密度大}，体积小\quad（$f$小且$m$不太大）
\item 线：轻/\unclear{刚}/\unclear{优}/\unclear{平}/长、（超过$1\unit{m}$）
\item 夹子：固定悬点，调节摆长
\item 秒表：教材实验中\ \fbox{仅有的使用秒表实验}
\end{itemize}
%%%END
%%%BLOCK id=108-02 ch=08 sec=单摆实验 kind=knowledge exp=1 alt=none cont=none y=0.17-0.22
\textbf{②}\ 摆长\quad $l=\dfrac{d}{2}+l_0$
%FIG p108_a 0.312 0.181 0.408 0.218
\notefig[0.25]{figures/p108_a.png}
摆长$=$半径$+$绳长.
%%%END
%%%BLOCK id=108-03 ch=08 sec=等效单摆 kind=knowledge alt=none cont=none y=0.22-0.30
\textbf{③}\ 等效单摆
%FIG p108_b 0.158 0.226 0.248 0.293
\notefig[0.25]{figures/p108_b.png}
$l=R\sin\theta$
%FIG p108_c 0.342 0.226 0.43 0.287
\notefig[0.25]{figures/p108_c.png}
$l=R-r$
%%%END
%%%BLOCK id=108-04 ch=08 sec=单摆实验 kind=method exp=1 alt=none cont=none y=0.30-0.35
\textbf{④}\ 周期：
\begin{itemize}
\item $n$次全振动用时$t$　$T=\dfrac{t}{n}$　\underline{从通过平衡位移开始计次数} % CHECK: 原稿此处写“位移”，疑为“位置”
\item $n$次过平衡位置　$T=\dfrac{2t}{n-1}$　　坑：从$0$开始数　$T=\dfrac{2t}{n}$
\end{itemize}
%%%END
%%%BLOCK id=108-05 ch=08 sec=等效重力加速度 kind=formula alt=02,03 cont=none y=0.35-0.42
\textbf{⑤}\ $g$（是等效$g$）
%FIG p108_d 0.243 0.368 0.306 0.418
\notefig[0.25]{figures/p108_d.png}
\[
\Rightarrow\ g'=\sqrt{g^2+a^2}
\]
%FIG p108_e 0.466 0.368 0.549 0.416
\notefig[0.25]{figures/p108_e.png}
\[
a=\sqrt{a_1^2+a_2^2+2a_1a_2}
\] % CHECK: 图中a1、a2之间标有夹角θ，式中疑缺cosθ，照原样转录
%%%END
%%%BLOCK id=108-06 ch=08 sec=单摆周期公式 kind=formula alt=14 cont=none y=0.42-0.49
\textbf{⑥}
\[
T=2\pi\sqrt{\dfrac{L}{g}}\qquad T^2=\dfrac{4\pi^2L}{g}\qquad \boxed{g=\dfrac{4\pi^2L}{T^2}}
\]
在圆弧运动时间$=$单摆周期　、弹簧\unclear{弹}的时间用$T=2\pi\sqrt{\dfrac{m}{k}}$　、LC电路充放电时间$T=2\pi\sqrt{LC}$
%%%END
%%%BLOCK id=108-07 ch=08 sec=单摆 kind=knowledge alt=03 cont=none y=0.48-0.57
%FIG p108_f 0.08 0.49 0.2 0.566
\notefig[0.25]{figures/p108_f.png}
同时放会相遇在平衡处\\
$L$　$g$一样　与释放角度无关，
%FIG p108_g 0.487 0.488 0.607 0.562
\notefig[0.25]{figures/p108_g.png}
单摆合力不为$0$　力变加速运动（“加速”二字上方另写有“减速”）\\
单摆运动没有平衡点.
%%%END
%%%BLOCK id=108-08 ch=08 sec=单摆实验 kind=method exp=1 alt=none cont=none y=0.57-0.79
\textbf{图像变换，}
%FIG p108_h 0.056 0.612 0.26 0.708
\notefig[0.3]{figures/p108_h.png}
\[
T^2=\dfrac{4\pi^2}{g}L
\]
$\pi^2=9.86$\\
改变摆长求$g$　$\big|$　用斜率求$g$
%FIG p108_i 0.308 0.598 0.562 0.728
\notefig[0.3]{figures/p108_i.png}
\[
T=2\pi\sqrt{\dfrac{L+R}{g}}
\]
\[
T^2=\dfrac{4\pi^2L}{g}+\dfrac{4\pi^2R}{g}
\]
%FIG p108_j 0.54 0.598 0.822 0.762
\notefig[0.35]{figures/p108_j.png}
\[
T=2\pi\sqrt{\dfrac{L+2R}{g}-\dfrac{R}{g}}
\]
\[
T^2=\dfrac{4\pi^2}{g}(L+2R)-\dfrac{4\pi^2}{g}R
\]
用纵截距求$k$ % CHECK: 末字像小写k，疑为R（由纵截距求球半径R），存疑
%%%END
%%%BLOCK id=108-09 ch=08 sec=单摆实验 kind=method exp=1 alt=none cont=none y=0.79-0.97
\fbox{两摆法}，剪$\Delta L$，不用测摆长
%FIG p108_k 0.02 0.866 0.145 0.932
\notefig[0.25]{figures/p108_k.png}
\[
T_1=2\pi\sqrt{\dfrac{L+\Delta L}{g}}\quad\text{①}
\]
\[
T_2=2\pi\sqrt{\dfrac{L}{g}}\quad\text{②}
\]
\[
\text{①}^2-\text{②}^2\qquad g=\dfrac{4\pi^2\Delta L}{T_1^2-T_2^2}
\]
%%%END
%%%BLOCK id=108-10 ch=08 sec=单摆实验 kind=method exp=1 alt=none cont=none y=0.85-0.97
单摆作打点计时器
%FIG p108_l 0.457 0.858 0.592 0.968
\notefig[0.25]{figures/p108_l.png}
%FIG p108_m 0.6 0.882 0.76 0.927
\notefig[0.3]{figures/p108_m.png}
相邻两点　时间间隔\\
$t=\pi\sqrt{\dfrac{L}{g}}=\dfrac{T}{2}$
%%%END
%%%BLOCK id=109-01 ch=08 sec=简谐运动·回复力 kind=knowledge alt=none cont=none y=0.03-0.15
\fbox{简谐运动}\quad $F_{\text{回}}=-kx$　判别式\\[0.8ex]
$a=-\dfrac{kx}{m}=\dfrac{\dd v}{\dd t}$　　$v=\dfrac{\dd x}{\dd t}$　　　振幅与机械能正相关\\[0.8ex]
回复力指向平衡位置（可变方向的\unclear{力}　　变力　不是不变、
%%%END
%%%BLOCK id=109-02 ch=08 sec=简谐运动·弹簧振子 kind=diagram alt=none cont=none y=0.10-0.40
\textbf{①}
%FIG p109_a 0.056 0.112 0.322 0.262
\notefig[0.4]{figures/p109_a.png}
% 图中标注：高 a_max v=0；平 v_max；低 a_max v=0（图中已含）
\[
T=2\pi\sqrt{\dfrac{m}{k}}
\]
%FIG p109_b 0.062 0.287 0.16 0.378
\notefig[0.25]{figures/p109_b.png}
%%%END
%%%BLOCK id=109-03 ch=08 sec=简谐运动·回复力证明 kind=derivation alt=02 cont=none y=0.12-0.40
\textbf{证明}
%FIG p109_c 0.455 0.147 0.646 0.272
\notefig[0.3]{figures/p109_c.png}
\begin{align*}
F_{\text{回}}&=mg+k\Delta x\\
\text{平衡时}\quad mg&=k(x-\Delta x)\\
\Rightarrow F_{\text{回}}&=kx\quad\text{又}\ F_{\text{回}}\text{与}x\text{反向}\\
\text{故}\quad F_{\text{回}}&=-kx
\end{align*}
%FIG p109_d 0.472 0.272 0.646 0.372
\notefig[0.3]{figures/p109_d.png}
\begin{align*}
F_{\text{回}}&=k\Delta x-mg\\
\text{平衡时}\quad mg&=k(\Delta x-x)\\
\therefore F_{\text{回}}&=-kx
\end{align*}
%%%END
%%%BLOCK id=109-04 ch=08 sec=简谐运动·回复力证明 kind=derivation alt=02 cont=none y=0.40-0.62
\textbf{②}　原长一致
%FIG p109_e 0.015 0.405 0.225 0.62
\notefig[0.3]{figures/p109_e.png}
\begin{align*}
F_{\text{回}}&=mg+k_1\Delta x_1+k_2\Delta x_1\\
mg&=k_1(x-\Delta x_1)+k_2(x-\Delta x_1)\\
\therefore F_{\text{回}}&=(k_1+k_2)x
\end{align*}
\underline{又因为$F_{\text{回}}$与$x$方向相反}\\
\struck{又反向}　故$F_{\text{回}}=-(k_1+k_2)x$
%%%END
%%%BLOCK id=109-05 ch=08 sec=简谐运动·对称性 kind=note alt=none cont=none y=0.40-0.46
对称　\fbox{等位移}\ \fbox{等速度}\ \fbox{等加速度}
%%%END
%%%BLOCK id=109-06 ch=08 sec=简谐运动·弹簧振子 kind=derivation alt=06 cont=none y=0.47-0.62
%FIG p109_f 0.645 0.474 0.835 0.585
\notefig[0.25]{figures/p109_f.png}
\begin{align*}
mgh&=\tfrac{1}{2}kh^2\\
mg&=k\Delta x\\
\therefore\Delta x&=\dfrac{h}{2}\\
\therefore h&=2\Delta x
\end{align*}
%%%END
%%%BLOCK id=109-07 ch=08 sec=简谐运动·浮体振动 kind=derivation alt=02 cont=none y=0.63-0.75
\textbf{③}
%FIG p109_g 0.066 0.658 0.185 0.745
\notefig[0.25]{figures/p109_g.png}
\[
F_{\text{回}}=G-F_{\text{浮}}=\rho gS(\Delta x+x)-\rho gS\Delta x=\rho gSx=kx\quad\text{又}\ F\ x\text{反向故}
\] % CHECK: 原稿“G-F浮=ρgS(Δx+x)-ρgSΔx”两项次序与G-F浮的符号不符(应为ρgSΔx-ρgS(Δx+x))，照原样转录
\[
F=-kx
\]
%%%END
%%%BLOCK id=109-08 ch=08 sec=单摆 kind=derivation alt=06 cont=none y=0.75-0.93
\textbf{④}　单摆的力　功　能　位移　加速度、周期分析.
%FIG p109_h 0.024 0.77 0.19 0.925
\notefig[0.3]{figures/p109_h.png}
% 图中标注：θ、a径、Ep max / Ek min、Ek max / Ep min（图中已含）
\[
F_{\text{回}}=mg\sin\theta=mg\dfrac{\widehat{l}}{r}=mg\dfrac{x}{r}=\dfrac{mg}{r}x=kx
\]
$\theta$很小时　$\theta\approx\tan\theta\approx\sin\theta$\\[0.8ex]
振幅$A$与$\bar v_{\text{下}}$，和$T$成正比　　$A=\bar v\dfrac{T}{4}$　　$T=mg\cos\theta+\dfrac{mv^2}{L}$\\[0.8ex]
机械能守恒　　$a_{\text{切}}=g\sin\theta$　　$a_{\text{径}}=\dfrac{V_{\unclear{\mathrm{max}}}^2}{L}=\dfrac{v^2}{L}$
%%%END
%%%BLOCK id=110-01 ch=08 sec=弹簧振子·弹簧模型题 kind=problem alt=06 cont=none y=0.03-0.13
\begin{problem}[①]
%FIG p110_a 0.105 0.04 0.24 0.125
\notefig[0.25]{figures/p110_a.png}
始终不分离　求最大压缩量
\begin{solution}
\[
\Delta X=\dfrac{2(M+m)g}{k}
\]
\end{solution}
\end{problem}
%%%END
%%%BLOCK id=110-02 ch=08 sec=弹簧振子·弹簧模型题 kind=problem alt=02 cont=none y=0.13-0.31
\begin{problem}[② 离地模型]
%FIG p110_b 0.088 0.182 0.222 0.31
\notefig[0.25]{figures/p110_b.png}
$B$恰好离地　求$A$向上距离
\begin{solution}
\begin{align*}
\Delta x_1&=\dfrac{m_Ag}{k}\\
\Delta x_2&=\dfrac{m_Bg}{k}\\
S=\Delta x_1+\Delta x_2&=\dfrac{(m_A+m_B)g}{k}
\end{align*}
\end{solution}
\end{problem}
%%%END
%%%BLOCK id=110-03 ch=08 sec=弹簧振子·弹簧模型题 kind=problem alt=02 cont=none y=0.31-0.45
\begin{problem}[③]
%FIG p110_c 0.057 0.34 0.24 0.45
\notefig[0.25]{figures/p110_c.png}
用$F$压一下后，$B$刚好被提起　求$F$
\begin{solution}
\[
F+3mg=k\left(\dfrac{4mg}{k}+\dfrac{3mg}{k}\right)
\]
\[
\therefore F=4mg
\]
\end{solution}
\end{problem}
%%%END
%%%BLOCK id=110-04 ch=08 sec=弹簧振子·弹簧模型题 kind=problem alt=03 cont=none y=0.46-0.57
\begin{problem}[④]
%FIG p110_d 0.108 0.462 0.25 0.558
\notefig[0.25]{figures/p110_d.png}
球在最高点时$M$恰好离开地面　求$M$加速度 % CHECK: 题末写“求M加速度”，而下面解出的a是球m的加速度，疑应为“求m加速度”
\begin{solution}
\[
(M+m)g=M\cdot 0+ma\qquad\text{\unclear{系}统牛二}
\]
\[
a=\left(\dfrac{M+m}{M}\right)g
\]% CHECK: 分母写作M；由(M+m)g=ma应为m，照原样转录
\end{solution}
\end{problem}
%%%END
%%%BLOCK id=110-05 ch=08 sec=弹簧振子·弹簧模型题 kind=problem alt=06,07 cont=none y=0.58-0.93
\begin{problem}[⑤]
%FIG p110_e 0.06 0.588 0.262 0.798
\notefig[0.25]{figures/p110_e.png}
图中标注：$C$　$m$；$h$；$A$　$m$　\unclear{未}粘性；$B$　$m$\\
$AC$到最高点时$B$与地面无压力
\begin{solution}
\[
V_1=\sqrt{2gh}\qquad mv=2mV_{\text{共}}=\dfrac{\sqrt{2gh}}{2}
\]% CHECK: 原稿为连等写法，实际应为V共=√(2gh)/2
$AC$在最高点时
\[
3mg=2ma\qquad\therefore a=\dfrac{3}{2}g
\]
\[
\left(\begin{aligned}
&2mg\cdot2\Delta x=\dfrac{1}{2}mV_{\text{共}}^2\qquad k=\dfrac{8mg}{h}\\
&\Delta x=\dfrac{mg}{k}
\end{aligned}\right)
\]% CHECK: 原稿动能项写作½mV共²；要得到k=8mg/h应为½(2m)V共²，照原样转录
\struck{$kx-2mg=2m\cdot\dfrac{3}{2}g$}　% 原稿此式整行被划去

\textbf{从功能角度}
\[
\text{OR}\quad\dfrac{1}{2}\cdot2m\left(\dfrac{\sqrt{2gh}}{2}\right)^2=2mg\Delta x+\dfrac{1}{2}k\left(\dfrac{2mg}{k}\right)^2
\]
\[
\Delta x=\dfrac{mg}{k}
\]
\end{solution}
\end{problem}
%%%END
%%%BLOCK id=111-01 ch=08 sec=简谐运动·图像 kind=diagram alt=none cont=none y=0.0-0.13
\unclear{手段大子}　% 页面顶端被裁切，字迹不全
%FIG p111_a 0.012 0.0 0.77 0.127
\notefig[0.85]{figures/p111_a.png}
% 图中标注：四幅图依次为(纵轴记号不清)、v、a、E(Ek、Ep)对t的曲线，图上标有T/2、T、t（图中已含）
%%%END
%%%BLOCK id=192-10 ch=08 sec=波的干涉·页边草算 kind=derivation alt=none cont=none y=0.32-0.70
% 页面右侧空白处的草算/草图(无题干；疑为两波源干涉：λ=20、T=1.5，“减”疑指减弱)；另在左下角(四星系统下方)有一幅无标注的曲线草图，疑与之有关
%FIG p192_d 0.71 0.325 0.985 0.445
\notefig[0.3]{figures/p192_d.png}
%FIG p192_c 0.02 0.595 0.145 0.697
\notefig[0.25]{figures/p192_c.png}

$\lambda=20$ % 写在上图下方偏右

$60-2x=10$\ \unclear{20m}

$2x-60=0$

\begin{itemize}
\item 竖排一列(左侧标\unclear{$L$})：10，20，40
\item 竖排一列(右侧)：10，20，40，60
\item 大括号内：$x=30$，$x=10$，$x=0$
\item 另有零散：$x=\unclear{10}$，$x=20$，$x=\unclear{25}$ % CHECK: 页边草算，排布及数字难以还原
\end{itemize}
%%%END
%%%BLOCK id=229-02 ch=08 sec=简谐运动·弹簧振子 kind=problem alt=03 cont=none y=0.46-0.76
\textbf{2.}
%FIG p229_b 0.01 0.455 0.33 0.60
\notefig[0.3]{figures/p229_b.png}

\begin{problem}
将C下压后释放，C作简谐运动，当C在最高点时，B压地

$k=100\unit{N/m}$ \quad 压力为0。求C到最低点时，C的加速度及$N_{\text{地}B}$
\begin{solution}
刚开始 \quad $kx_0=mg$ \quad $x_0=0.02\unit{m}$

C向下$x_1$，振幅 $A=x_1$

C在最高处 \quad $Mg=k(A-x_0)$ \quad $\therefore A=0.07\unit{m}$

C在最低处 \quad $k(A+x_0)-mg=ma$ \quad $a=35\unit{m/s^2}$

对B \quad $F=mg+k(A+x_0)=14\unit{N}$ % CHECK: 按数值(5+9=14)此处应为Mg
\end{solution}
\end{problem}
%%%END
%%%BLOCK id=232-02 ch=08 sec=波的叠加·多解与最小公倍数 kind=problem alt=none cont=none y=0.12-0.235
\begin{problem}[2]
\begin{solution}
$x_{\text{甲}}=50+50n_1$ \qquad $V_{\text{甲}}=V_{\text{乙}}=25\unit{cm/s}$\par
$x_{\text{乙}}=50+60n_2$ $\;\to\;$ $x=50+300n$ \quad 找最小公倍数\par
波程差 $\Delta x=160n_2-50n_{\text{甲}}+5\ge 5$ % CHECK: 160 的首位"1"疑为多余笔画（按 x 乙 的式子似应为 60n_2）
\qquad $t_{\min}=\dfrac{\Delta x}{\text{\unclear{$\blacksquare$}}V_{\text{相}}}=0.15$
\end{solution}
\end{problem}
%%%END
%%%BLOCK id=232-03 ch=08 sec=波的干涉·减弱点 kind=problem alt=none cont=none y=0.25-0.33
\begin{problem}[3]
%FIG p232_b 0.07 0.262 0.20 0.30
\notefig[0.2]{figures/p232_b.png}
\begin{solution}
设$x$为\unclear{减弱}点
\[ x-(2-x)=(2n+1)\frac{\lambda}{2} \qquad x=\frac{1}{2}n+\frac{5}{4} \qquad x\in[0,2] \]
故 $n=-2,-1,0,1$
\end{solution}
\end{problem}
%%%END
%%%BLOCK id=274-01 ch=08 sec=波动图像·起振方向 kind=problem alt=none cont=none y=0.00-0.29
% 页首为贴上的印刷体题目(页顶/页右被截断，仅见部分题干)，图上与图旁有手写批注
\begin{problem}
\unclear{…}简谐运动，$t_1=3\unit{s}$ 时，波刚传到 $x=6\unit{m}$ 处，波形如图中实\unclear{线所}示，经过时间 $t_2$ 形成的部分波形如图虚线所示，则波源的起振方向为\underline{\ 沿$y$轴\unclear{正}向\ }\unclear{…}$t_2$ 的最短时间为\underline{\ 5.5\ }s。 % 两处下划线内为手写填空

%FIG p274_a 0.37 0.055 0.73 0.212
\notefig[0.5]{figures/p274_a.png}
\begin{solution}
$\lambda=4\unit{m}$

$v=\dfrac{6}{3}=2\unit{m/s}$

到这才是起振方向。 % 手写箭头由此句指向图中 $x$ 轴右端的波前处

$t_2=\dfrac{11}{2}=5.5\unit{s}$
\end{solution}
\end{problem}
%%%END
%%%BLOCK id=281-06 ch=08 sec=波动图像与振动图像·多解问题 kind=problem alt=none cont=none y=0.53-1.00
% 贴纸题（印刷体）+手写解答。原稿：“（10分）”处有一道斜划线；“波长大于9cm，”与“1cm，”被圈出，“且传播时无衰减”下有横线；“运动，”与“t=2s”之间有一道斜划线；①中“B”下有弧线
\begin{problem}[10分]
均匀介质中质点$A$、$B$的平衡位置位于$x$轴上，坐标分别为$x_A=-1\unit{cm}$和$x_B=8\unit{cm}$。某简谐横波沿$x$轴正方向传播，波长大于$9\unit{cm}$，振幅为$1\unit{cm}$，且传播时无衰减。$t=0$时刻开始计时，质点$A$位于$x$轴上方$0.5\unit{cm}$处，质点$B$位于平衡位置且向$x$轴下方运动，$t=2\unit{s}$质点$A$第一次回到它$t=0$时刻的位置。求：

①从$t=0$时刻开始，质点$B$经过多长时间第一次位于波峰；

②该简谐波的波速。
\begin{solution}
设$A$质点振动方程\quad $y_A=A\sin\Bigl(\dfrac{2\pi}{T}t+\varphi\Bigr)$

$t=0$时$A$对应初相为$\dfrac{\pi}{6}$，or $\dfrac{5\pi}{6}$

%FIG p281_d 0.04 0.815 0.30 0.91
\notefig[0.3]{figures/p281_d.png}
②\ $t=2$时\quad $\dfrac{2\pi}{T_2}t+\dfrac{5\pi}{6}=\dfrac{13\pi}{6}$

$T_2=3\unit{s}$

经$\dfrac{3}{4}T_2=2.25$第一次位于波峰

①\ $t=2\unit{s}$时\quad $\dfrac{2\pi}{T_1}t+\dfrac{\pi}{6}=\dfrac{5\pi}{6}$

$T_1=6\unit{s}$

经$\dfrac{3}{4}T=4.15\unit{s}$第一次位于波峰% CHECK: 手写为“4.15”（按 3/4×6=4.5 应为 4.5s，疑多写了“1”）；此处 T 未见下标

②\ [①]\ $\Delta x=x_B-x_A=\dfrac{7}{12}\lambda$

$\lambda_1=\dfrac{108}{7}\unit{cm}$

[②]\ $\Delta x=x_B-x_A=\dfrac{11}{12}\lambda$

$\lambda_2=\dfrac{108}{11}\unit{cm}$

$v=\dfrac{\lambda}{T}$

$v_1=\dfrac{18}{7}\unit{cm/s}$

or $v_2=\dfrac{36}{11}\unit{cm/s}$
\end{solution}
\end{problem}
%%%END
%%%BLOCK id=283-01 ch=08 sec=波的叠加·波峰重合 kind=problem alt=none cont=none y=0.00-0.43
% 题干为粘贴的印刷题条（粘贴条上下行有错位，此处按语义顺序转录）；手写在“中沿x轴传播”“轴正向”“乙沿x轴负向传播”“二者振动方向一致”下画了线
（图上方手写标题）\unclear{常规解法}

%FIG p283_a 0.545 0.0 0.95 0.105
\notefig[0.5]{figures/p283_a.png}

\begin{problem}[19.（6分）]
甲、乙两列简谐横波在同一介质中传播的波速均为 $v=2\unit{m/s}$。当甲单独传播时，在空间中某质点引起的振动位移随时间变化关系如图（a）所示，当乙单独传播时，在空间中某质点引起的振动位移随时间化关系如图（b）所示。% 印刷原文此处漏“变”字
现使两列波同时在该介质中沿 $x$ 轴传播，甲沿 $x$ 轴正向、乙沿 $x$ 轴负向传播，二者振动方向一致，在 $t=0$ 时刻，平衡位置坐标 $x=8\unit{cm}$ 的点正好位于偏离平衡位置 $+8\unit{cm}$ 处。

（1）$t=0$ 时，求介质中偏离平衡位置位移为 $+8\unit{cm}$ 的所有质点的 $x$ 坐标

（2）求 $x=8\unit{cm}$ 的点下一次出现偏离平衡位置 $+8\unit{cm}$ 的时间 $t_1$ 及之后介质中有质点出现偏离平衡位置 $+8\unit{cm}$ 的最短时间 $t_2$

\begin{solution}
[\unclear{找公倍数}]

（1）甲\quad $T_1=0.4\unit{s}$\quad $\lambda_1=vT_1=80\unit{cm}$\quad 甲波峰 $x=8+80m\ (\unit{cm})$\quad $(m=0,\pm1,\pm2,\cdots)$

\hspace{1.6em}乙\quad $T_2=0.5\unit{s}$\quad $\lambda_2=vT_2=100\unit{cm}$\quad 乙波峰 $x=8+100n\ (\unit{cm})$\quad $(n=0,\pm1,\pm2,\cdots)$

\hspace{1.6em}$\therefore$ 甲乙同时波峰对应 $x=8+400k\ (\unit{cm})$\quad $(k=0,\pm1,\pm2,\cdots)$

（2）甲波峰每隔 $0.4\unit{s}$ 到达 $x=8\unit{cm}$，乙波峰每隔 $0.5\unit{s}$ 到该处

\hspace{1.6em}故 $x=8\unit{cm}$ 点下一次出现同时波峰用时 $t_1=2\unit{s}$

\hspace{1.6em}$t=0$ 时甲乙波峰除了重叠处以外，最近的波峰相距 $20\unit{cm}$

\hspace{1.6em}$\therefore$ 下一次有质点出现偏离平衡位置 $+8\unit{cm}$ 的时间为 $t_2=\dfrac{\Delta x}{2v}=0.05\unit{s}$
\end{solution}
\end{problem}
%%%END
%%%BLOCK id=289-03 ch=08 sec=机械振动·质点位移相同 kind=note alt=none cont=none y=0.28-0.37
\unclear{无搅3算}\quad $x^2=g$% CHECK: 首行第2、4字不清（第4字被反复涂写）；“x”也可能是 $\omega$

每隔 $t$ 两\unclear{个}质点位移相同一次\quad $2t=nT$
%%%END
